\section{Introduction}
\label{sec:introduction}

Resonances of area-preserving and reversible maps provide a natural setting in which local bifurcation theory, Hamiltonian normal forms, and periodic-orbit geometry interact. Among the simplest nontrivial cases is the \(1{:}2\) resonance, characterized by a double multiplier \(-1\) of the Poincaré return. At the resonant parameter, the square of the return is tangent to the identity, so that nonlinear terms determine the local organization of nearby periodic orbits.

Reversibility adds further structure. Symmetric periodic orbits intersect fixed sets of reversing involutions, and the corresponding symmetry lines provide a natural framework for locating and classifying bifurcating branches \cite{Sevryuk1986,LambRoberts1998}. In the area-preserving setting, normal forms near neutral multipliers \(+1\) and \(-1\) were studied systematically by Gelfreich and Gelfreikh \cite{GelfreichGelfreikh2010}. Near a semisimple multiplier \(-1\), the return admits a formal odd normal form, so that after changing sign one obtains a near-identity exact symplectic map described by an even interpolating Hamiltonian. Interpolating Hamiltonians have long been used to study resonant geometry and stability in conservative maps \cite{BazzaniEtAl1993}.

The system considered here arises from the normal geodesic flow of a rank-two, seven-dimensional Carnot structure with growth vector
\[
(2,3,5,7).
\]
The corresponding nilpotent algebra and its sub-Riemannian geodesic problem have previously appeared in the literature. In particular, Kruglikov, Vollmer, and Lukes-Gerakopoulos \cite{KruglikovVollmerLukesGerakopoulos2017} study the same structure, its coadjoint invariants, and questions of polynomial integrability. General background on sub-Riemannian extremals, left-invariant Hamiltonian systems, and coadjoint reduction may be found in \cite{AgrachevSachkov2004,AgrachevBarilariBoscain2019}.

Our starting point is therefore not a new Carnot algebra. The contribution begins with the reduced dynamics on its generic coadjoint orbits and, specifically, with a resonant periodic family that has not previously been analyzed in this form.

After rotational reduction, Carnot scaling, and restriction to unit speed, the geodesic equations reduce to
\begin{equation}
\boxed{
\dot x=\cos z,\qquad
\dot y=\sin z,\qquad
\dot z=-\frac12x^2+\frac12y^2-\mu.
}
\label{eq:intro-reduced-flow}
\end{equation}
The parameter \(\mu\) is the scale-invariant combination of the coadjoint Casimirs. The flow preserves volume and induces an exact symplectic local Poincaré map on the section \(x=0\).

Continuation of a primary symmetric periodic orbit numerically locates a \(1{:}2\) resonance at
\[
\mu=\mu_\star\simeq0.662869.
\]
At the corresponding parameter, the full monodromy matrix is numerically consistent with
\[
DP_{\mu_\star}=-I.
\]

The distinctive feature appears after fixing the linear canonical normalization by two pieces of information: the infinitesimal detuning and the tangent fixed sets of the two reversing involutions. In this fixed reversible chart the leading nonlinear Hamiltonian is
\[
H_4=A(Q^4+P^4)+BQ^2P^2.
\]
Writing
\[
\rho^2=2-\frac BA,
\]
we find
\[
\rho\simeq2.00118.
\]
The exactly quartic-degenerate value is \(\rho=2\), for which
\[
H_4=A(Q^2-P^2)^2.
\]

Thus the present system lies extremely close to a quartic degeneracy in which the diagonal directions are null directions of the leading nonlinear Hamiltonian. Define
\[
\sigma=4-\rho^2.
\]
Numerically,
\[
\sigma\simeq-4.73\times10^{-3}.
\]
The smallness of this coefficient promotes the sixth-order Hamiltonian
\[
H_6=dQ^6+eQ^4P^2+fQ^2P^4+gP^6
\]
to leading order in the near-diagonal sector.

The result is a two-scale local geometry. Along the axial directions,
\[
I_H\asymp\delta,\qquad |\mcJ_H|\asymp\delta^2.
\]
Near the almost-null diagonal directions, the leading scalar reduction is
\[
K_E(I)=-\delta I+A\sigma I^2+GI^3,
\qquad G=d+e+f+g>0,
\]
and in the regime
\[
\deltax\ll\delta\ll1
\]
one obtains
\[
I_E\asymp\delta^{1/2},
\qquad
|\mcJ_E|\asymp\delta^{3/2}.
\]

The \(3/2\) law is an intermediate asymptotic in the sense of scaling regimes lying between microscopic and ultimate asymptotic behavior \cite{Barenblatt1996}. Since \(\sigma<0\) is fixed, the branch ultimately crosses over to a finite-amplitude plateau as \(\delta\to0^+\). Direct numerical solution of the exact second-return equation at resonance reveals the corresponding residual elliptic period-two cycles.

The phenomenon considered here is distinct from the renormalization theory of repeated conservative period doubling \cite{GaidashevKoch2011}. We study a single semisimple local \(1{:}2\) resonance and the promotion of the sixth-order normal form caused by proximity to a quartic degeneracy; no period-doubling cascade or renormalization fixed point is assumed.

The novelty of the present work is therefore not the \((2,3,5,7)\) Carnot algebra itself, but the identification and analysis of a near-degenerate reversible \(1{:}2\) resonance in its reduced geodesic flow. The quartic near-degeneracy promotes the sixth-order normal form and produces two distinct periodic-orbit action scales together with a crossover to residual elliptic period-two cycles.

The exactly quartic-degenerate limit \(\sigma=0\) defines a separate higher-codimension problem and is not analyzed here.
